\section{Conclusion}
\label{sec:conclusion}

Judicial enforcement changes public procurement through more than prices. In S\~ao Paulo's urgent pharmaceutical markets, court-mandated delivery raises costs relative to the closest administrative urgent channel, but the same supplier does not broadly charge more to the same buyer for the same item in deep markets. The gap instead appears through smaller orders and different winning suppliers.

That distinction matters for policy. If the problem were only same-firm markups, price caps or contract clauses would be the natural response. If the problem is sourcing, the institutional response is different: preserve legal delivery while rebuilding aggregation, pre-contracting, and supplier matching under urgency. The evidence supports the latter interpretation for deep repeated urgent markets, while leaving room for supplier leverage in thinner ones.
